% Options for packages loaded elsewhere
\PassOptionsToPackage{unicode}{hyperref}
\PassOptionsToPackage{hyphens}{url}
\documentclass[
  french,
  10pt,
  a4paper,
]{article}
\usepackage{xcolor}
\usepackage[margin=1.8cm]{geometry}
\usepackage{amsmath,amssymb}
\setcounter{secnumdepth}{-\maxdimen} % remove section numbering
\usepackage{iftex}
\ifPDFTeX
  \usepackage[T1]{fontenc}
  \usepackage[utf8]{inputenc}
  \usepackage{textcomp} % provide euro and other symbols
\else % if luatex or xetex
  \usepackage{unicode-math} % this also loads fontspec
  \defaultfontfeatures{Scale=MatchLowercase}
  \defaultfontfeatures[\rmfamily]{Ligatures=TeX,Scale=1}
\fi
\usepackage{lmodern}
\ifPDFTeX\else
  % xetex/luatex font selection
\fi
% Use upquote if available, for straight quotes in verbatim environments
\IfFileExists{upquote.sty}{\usepackage{upquote}}{}
\IfFileExists{microtype.sty}{% use microtype if available
  \usepackage[]{microtype}
  \UseMicrotypeSet[protrusion]{basicmath} % disable protrusion for tt fonts
}{}
\makeatletter
\@ifundefined{KOMAClassName}{% if non-KOMA class
  \IfFileExists{parskip.sty}{%
    \usepackage{parskip}
  }{% else
    \setlength{\parindent}{0pt}
    \setlength{\parskip}{6pt plus 2pt minus 1pt}}
}{% if KOMA class
  \KOMAoptions{parskip=half}}
\makeatother
\usepackage{color}
\usepackage{fancyvrb}
\newcommand{\VerbBar}{|}
\newcommand{\VERB}{\Verb[commandchars=\\\{\}]}
\DefineVerbatimEnvironment{Highlighting}{Verbatim}{commandchars=\\\{\}}
% Add ',fontsize=\small' for more characters per line
\newenvironment{Shaded}{}{}
\newcommand{\AlertTok}[1]{\textcolor[rgb]{1.00,0.00,0.00}{\textbf{#1}}}
\newcommand{\AnnotationTok}[1]{\textcolor[rgb]{0.38,0.63,0.69}{\textbf{\textit{#1}}}}
\newcommand{\AttributeTok}[1]{\textcolor[rgb]{0.49,0.56,0.16}{#1}}
\newcommand{\BaseNTok}[1]{\textcolor[rgb]{0.25,0.63,0.44}{#1}}
\newcommand{\BuiltInTok}[1]{\textcolor[rgb]{0.00,0.50,0.00}{#1}}
\newcommand{\CharTok}[1]{\textcolor[rgb]{0.25,0.44,0.63}{#1}}
\newcommand{\CommentTok}[1]{\textcolor[rgb]{0.38,0.63,0.69}{\textit{#1}}}
\newcommand{\CommentVarTok}[1]{\textcolor[rgb]{0.38,0.63,0.69}{\textbf{\textit{#1}}}}
\newcommand{\ConstantTok}[1]{\textcolor[rgb]{0.53,0.00,0.00}{#1}}
\newcommand{\ControlFlowTok}[1]{\textcolor[rgb]{0.00,0.44,0.13}{\textbf{#1}}}
\newcommand{\DataTypeTok}[1]{\textcolor[rgb]{0.56,0.13,0.00}{#1}}
\newcommand{\DecValTok}[1]{\textcolor[rgb]{0.25,0.63,0.44}{#1}}
\newcommand{\DocumentationTok}[1]{\textcolor[rgb]{0.73,0.13,0.13}{\textit{#1}}}
\newcommand{\ErrorTok}[1]{\textcolor[rgb]{1.00,0.00,0.00}{\textbf{#1}}}
\newcommand{\ExtensionTok}[1]{#1}
\newcommand{\FloatTok}[1]{\textcolor[rgb]{0.25,0.63,0.44}{#1}}
\newcommand{\FunctionTok}[1]{\textcolor[rgb]{0.02,0.16,0.49}{#1}}
\newcommand{\ImportTok}[1]{\textcolor[rgb]{0.00,0.50,0.00}{\textbf{#1}}}
\newcommand{\InformationTok}[1]{\textcolor[rgb]{0.38,0.63,0.69}{\textbf{\textit{#1}}}}
\newcommand{\KeywordTok}[1]{\textcolor[rgb]{0.00,0.44,0.13}{\textbf{#1}}}
\newcommand{\NormalTok}[1]{#1}
\newcommand{\OperatorTok}[1]{\textcolor[rgb]{0.40,0.40,0.40}{#1}}
\newcommand{\OtherTok}[1]{\textcolor[rgb]{0.00,0.44,0.13}{#1}}
\newcommand{\PreprocessorTok}[1]{\textcolor[rgb]{0.74,0.48,0.00}{#1}}
\newcommand{\RegionMarkerTok}[1]{#1}
\newcommand{\SpecialCharTok}[1]{\textcolor[rgb]{0.25,0.44,0.63}{#1}}
\newcommand{\SpecialStringTok}[1]{\textcolor[rgb]{0.73,0.40,0.53}{#1}}
\newcommand{\StringTok}[1]{\textcolor[rgb]{0.25,0.44,0.63}{#1}}
\newcommand{\VariableTok}[1]{\textcolor[rgb]{0.10,0.09,0.49}{#1}}
\newcommand{\VerbatimStringTok}[1]{\textcolor[rgb]{0.25,0.44,0.63}{#1}}
\newcommand{\WarningTok}[1]{\textcolor[rgb]{0.38,0.63,0.69}{\textbf{\textit{#1}}}}
\usepackage{longtable,booktabs,array}
\newcounter{none} % for unnumbered tables
\usepackage{calc} % for calculating minipage widths
% Correct order of tables after \paragraph or \subparagraph
\usepackage{etoolbox}
\makeatletter
\patchcmd\longtable{\par}{\if@noskipsec\mbox{}\fi\par}{}{}
\makeatother
% Allow footnotes in longtable head/foot
\IfFileExists{footnotehyper.sty}{\usepackage{footnotehyper}}{\usepackage{footnote}}
\makesavenoteenv{longtable}
\ifLuaTeX
\usepackage[bidi=basic,shorthands=off]{babel}
\else
\usepackage[bidi=default,shorthands=off]{babel}
\fi
\ifLuaTeX
  \usepackage{selnolig} % disable illegal ligatures
\fi
\setlength{\emergencystretch}{3em} % prevent overfull lines
\providecommand{\tightlist}{%
  \setlength{\itemsep}{0pt}\setlength{\parskip}{0pt}}
\usepackage{bookmark}
\IfFileExists{xurl.sty}{\usepackage{xurl}}{} % add URL line breaks if available
\urlstyle{same}
\hypersetup{
  pdflang={fr},
  hidelinks,
  pdfcreator={LaTeX via pandoc}}

\author{}
\date{}

\begin{document}

\section{TP --- Sécuriser le backend d'un gestionnaire de
tâches}\label{tp-suxe9curiser-le-backend-dun-gestionnaire-de-tuxe2ches}

\textbf{Spring Boot, Spring Security et client React Planora}

Durée : \textbf{60 minutes}, après une présentation théorique de 30
minutes. Travail individuel ou en binôme. Technologies : Java 17, Maven,
Spring Boot et React/Vite.

\textbf{Nom et prénom :}
\ldots\ldots\ldots\ldots\ldots\ldots\ldots\ldots\ldots\ldots\ldots\ldots\ldots\ldots\ldots\ldots\ldots\ldots\ldots\ldots.\\
\textbf{Groupe :}
\ldots\ldots\ldots\ldots\ldots\ldots\ldots\ldots\ldots. \textbf{Date :}
\ldots\ldots\ldots\ldots\ldots\ldots\ldots\ldots\ldots.

\subsection{1. Le produit et le travail
demandé}\label{le-produit-et-le-travail-demanduxe9}

Planora est une application de gestion des tâches d'une équipe. Son
interface React et son API métier sont déjà fournies : on peut
consulter, créer, modifier, terminer et supprimer des tâches. Les
statistiques proviennent des tâches réellement stockées par Spring.

Le TP fait évoluer \textbf{le backend Spring}, en ajoutant
progressivement l'authentification puis les autorisations. Utilisez
l'application pour constater les effets de vos changements. Les règles
d'accès doivent être imposées par Spring ; aucune modification de React
n'est demandée.

À la fin, vous saurez distinguer authentification et autorisation,
déclarer des comptes avec BCrypt, configurer
\texttt{SecurityFilterChain}, protéger une méthode avec
\texttt{@PreAuthorize} et vérifier les parcours métier avec les tests et
les réponses HTTP.

Les tâches sont partagées par l'équipe : \textbf{Alice et Bob peuvent
gérer toutes les tâches}. Le TP ne met pas en place de propriété
individuelle des tâches. Les statistiques et le projet confidentiel sont
réservés à Bob, administrateur.

\subsubsection{Contrat de l'API}\label{contrat-de-lapi}

{\def\LTcaptype{none} % do not increment counter
\begin{longtable}[]{@{}
  >{\raggedright\arraybackslash}p{(\linewidth - 4\tabcolsep) * \real{0.3333}}
  >{\raggedright\arraybackslash}p{(\linewidth - 4\tabcolsep) * \real{0.3333}}
  >{\raggedright\arraybackslash}p{(\linewidth - 4\tabcolsep) * \real{0.3333}}@{}}
\toprule\noalign{}
\begin{minipage}[b]{\linewidth}\raggedright
Méthode et route
\end{minipage} & \begin{minipage}[b]{\linewidth}\raggedright
Fonction
\end{minipage} & \begin{minipage}[b]{\linewidth}\raggedright
Succès
\end{minipage} \\
\midrule\noalign{}
\endhead
\bottomrule\noalign{}
\endlastfoot
\texttt{GET\ /api/public/info} & Information publique, à consulter avec
curl & \texttt{200}, texte \\
\texttt{GET\ /api/tasks} & Liste des tâches & \texttt{200}, tableau
JSON \\
\texttt{POST\ /api/tasks} & Création & \texttt{201}, tâche JSON \\
\texttt{PUT\ /api/tasks/\{id\}} & Modification & \texttt{200}, tâche
JSON \\
\texttt{PATCH\ /api/tasks/\{id\}/complete} & Passage à l'état terminé &
\texttt{200}, tâche JSON \\
\texttt{DELETE\ /api/tasks/\{id\}} & Suppression & \texttt{204}, sans
corps \\
\texttt{GET\ /api/admin/stats} & Statistiques calculées & \texttt{200},
objet JSON \\
\texttt{GET\ /api/tasks/confidential} & Note du projet Orion, fournie à
partir de la branche 5 & \texttt{200}, objet JSON \\
\end{longtable}
}

Le CRUD et les statistiques existent dès la branche 1. La route
confidentielle est absente des branches 1 à 4 ; elle est déjà
implémentée dans la branche 5, où il reste à la protéger.

Une tâche contient \texttt{id}, \texttt{title}, \texttt{description},
\texttt{status}, \texttt{priority}, \texttt{dueDate} et
\texttt{assignee}. Le formulaire React permet de modifier le titre, la
description, l'état, la priorité et l'échéance. La personne assignée est
affichée ; elle n'est pas éditable dans ce formulaire et le serveur la
conserve lors d'une modification.

États : \texttt{TODO}, \texttt{IN\_PROGRESS}, \texttt{DONE}. Priorités :
\texttt{LOW}, \texttt{MEDIUM}, \texttt{HIGH}. Le titre est obligatoire
et limité à 160 caractères ; la description à 1 000 caractères ;
l'échéance utilise le format \texttt{AAAA-MM-JJ}. Des données invalides
renvoient \texttt{400} et un identifiant de tâche inexistant renvoie
\texttt{404}, si les filtres de sécurité laissent passer la requête.

Les données restent \textbf{en mémoire}. Un redémarrage de Spring
restaure les trois tâches d'exemple. Au démarrage, les statistiques sont
: 3 tâches, 1 terminée, 2 en attente, 1 de priorité haute et 33 \%
d'achèvement. Les statistiques d'administration sont calculées sur
toutes les tâches ; les compteurs de l'écran des tâches utilisent la
liste chargée dans React.

\subsubsection{Architecture}\label{architecture}

{\def\LTcaptype{none} % do not increment counter
\begin{longtable}[]{@{}ll@{}}
\toprule\noalign{}
Partie & Responsabilité \\
\midrule\noalign{}
\endhead
\bottomrule\noalign{}
\endlastfoot
\texttt{frontend/src/main.jsx} & Écrans métier et formulaire de
connexion \\
\texttt{frontend/src/api.js} & Appels JSON et en-tête HTTP Basic \\
\texttt{controller/} & Routes HTTP et appels des services \\
\texttt{model/} & Objets échangés avec React \\
\texttt{service/TaskService.java} & CRUD et stockage en mémoire \\
\texttt{service/StatisticsService.java} & Indicateurs calculés \\
\texttt{config/UserConfig.java} & Comptes, rôles et encodeur BCrypt \\
\texttt{config/SecurityConfig.java} & Règles HTTP et sécurité des
méthodes \\
\end{longtable}
}

Les chemins Java sont relatifs à
\texttt{src/main/java/fr/tpspringsecurity/}. Les fichiers de
configuration apparaissent au fil des branches. La même interface React
est présente sur les six branches.

\subsection{2. Comprendre les branches}\label{comprendre-les-branches}

Chaque branche est \textbf{le point de départ de son exercice}, avec les
corrections précédentes. La branche 6 contient la solution finale. Les
numéros ne désignent donc pas des exercices déjà terminés.

{\def\LTcaptype{none} % do not increment counter
\begin{longtable}[]{@{}
  >{\raggedright\arraybackslash}p{(\linewidth - 6\tabcolsep) * \real{0.2500}}
  >{\raggedright\arraybackslash}p{(\linewidth - 6\tabcolsep) * \real{0.2500}}
  >{\raggedright\arraybackslash}p{(\linewidth - 6\tabcolsep) * \real{0.2500}}
  >{\raggedright\arraybackslash}p{(\linewidth - 6\tabcolsep) * \real{0.2500}}@{}}
\toprule\noalign{}
\begin{minipage}[b]{\linewidth}\raggedright
Branche
\end{minipage} & \begin{minipage}[b]{\linewidth}\raggedright
État initial
\end{minipage} & \begin{minipage}[b]{\linewidth}\raggedright
Travail demandé
\end{minipage} & \begin{minipage}[b]{\linewidth}\raggedright
Durée
\end{minipage} \\
\midrule\noalign{}
\endhead
\bottomrule\noalign{}
\endlastfoot
\texttt{Tp\_version1} & CRUD et statistiques ouverts & Auditer les accès
& 10 min \\
\texttt{Tp\_version2} & API encore ouverte, dépendance à ajouter &
Installer Spring Security & 10 min \\
\texttt{Tp\_version3} & Sécurité par défaut, compte généré & Déclarer
Alice et Bob & 15 min \\
\texttt{Tp\_version4} & Alice et Bob présents, politique par défaut &
Configurer les droits HTTP et CSRF & 10 min \\
\texttt{Tp\_version5} & Droits HTTP corrigés ; note lisible par tout
compte valide & Activer la protection de méthode & 10 min \\
\texttt{Tp\_version6} & Sécurité complète & Vérifier le parcours final &
5 min \\
\end{longtable}
}

Après l'exercice 2 et après l'exercice 3, les lectures demandent un
compte valide, y compris l'information publique. Les rôles ne limitent
pas encore les statistiques. \textbf{Sans jeton CSRF, les mutations
renvoient \texttt{403}, même avec un compte valide et même sans compte.}
L'exercice 4 adapte cette configuration pour permettre les mutations du
client fourni.

Le compte généré \texttt{user} apparaît après l'ajout de la dépendance
en exercice 2 et reste présent au début de la branche 3. Alice et Bob
existent \textbf{après la réalisation de l'exercice 3}, et dès
l'ouverture des branches 4 à 6.

\subsection{3. Préparer et lancer le
projet}\label{pruxe9parer-et-lancer-le-projet}

Prérequis : JDK 17, Maven 3.6.3 ou supérieur, Git et Node.js
\textbf{20.19+ ou 22.12+} avec npm. Vérifiez notamment que
\texttt{mvn\ -version} utilise le JDK attendu.

\begin{Shaded}
\begin{Highlighting}[]
\ExtensionTok{java} \AttributeTok{{-}version}
\ExtensionTok{mvn} \AttributeTok{{-}version}
\ExtensionTok{node} \AttributeTok{{-}{-}version}
\ExtensionTok{npm} \AttributeTok{{-}{-}version}
\end{Highlighting}
\end{Shaded}

Récupérez le dépôt :

\begin{Shaded}
\begin{Highlighting}[]
\FunctionTok{git}\NormalTok{ clone https://github.com/romarich57/Tp\_Spring\_security.git}
\BuiltInTok{cd}\NormalTok{ Tp\_Spring\_security}
\FunctionTok{git}\NormalTok{ switch Tp\_version1}
\end{Highlighting}
\end{Shaded}

Dans un premier terminal, à la racine :

\begin{Shaded}
\begin{Highlighting}[]
\ExtensionTok{mvn}\NormalTok{ test}
\ExtensionTok{mvn}\NormalTok{ spring{-}boot:run}
\end{Highlighting}
\end{Shaded}

Dans un second terminal :

\begin{Shaded}
\begin{Highlighting}[]
\BuiltInTok{cd}\NormalTok{ frontend}
\ExtensionTok{npm}\NormalTok{ install}
\ExtensionTok{npm}\NormalTok{ run dev}
\end{Highlighting}
\end{Shaded}

Ouvrez \textbf{http://localhost:5173}. L'API écoute sur
\textbf{http://localhost:8080}. Vite relaie \texttt{/api} vers Spring,
ce qui permet à React de travailler sur la même origine via le proxy.

\subsubsection{Utiliser Planora}\label{utiliser-planora}

\begin{itemize}
\tightlist
\item
  \textbf{Mes tâches} : consulter, rechercher, filtrer, créer, modifier,
  terminer et supprimer.
\item
  \textbf{Statistiques} : consulter les indicateurs du backend.
\item
  \textbf{Projet confidentiel} : lire la note lorsque le backend propose
  cette route et autorise l'accès.
\item
  \textbf{Connexion} : saisir le nom du compte et son mot de passe ; se
  déconnecter pour revenir au parcours anonyme.
\end{itemize}

Le formulaire transmet les identifiants en HTTP Basic à chaque requête.
Ils restent en mémoire React, sans cookie ni stockage local ; recharger
la page les efface. Il n'existe pas de route de login dédiée.
L'affichage du nom saisi ne prouve pas que le compte est valide : Spring
vérifie les identifiants à chaque appel et peut renvoyer \texttt{401}.

L'interface affiche les données et les erreurs réelles du serveur. Pour
relever précisément un statut de succès, utilisez l'onglet Réseau du
navigateur, \texttt{requests.http} ou \texttt{curl}. L'information
publique n'a pas d'écran dans Planora : vérifiez-la avec \texttt{curl}.

\subsubsection{Passer à la branche
suivante}\label{passer-uxe0-la-branche-suivante}

Arrêtez Spring et Vite avec \texttt{Ctrl+C}. Enregistrez les fichiers de
votre exercice dans un commit local, puis changez de branche et relancez
les deux serveurs. Exemple après l'exercice 2 :

\begin{Shaded}
\begin{Highlighting}[]
\FunctionTok{git}\NormalTok{ add pom.xml src/test/java/fr/tpspringsecurity/TaskApiSecurityTest.java}
\FunctionTok{git}\NormalTok{ commit }\AttributeTok{{-}m} \StringTok{"Exercice 2 terminé"}
\FunctionTok{git}\NormalTok{ switch Tp\_version3}
\end{Highlighting}
\end{Shaded}

Les branches suivantes fournissent déjà les corrections précédentes ; il
n'est pas nécessaire de reporter votre commit sur elles. Vous pouvez
aussi mettre vos changements de côté avec
\texttt{git\ stash\ push\ -m\ "Mon\ exercice"} ; ne réappliquez pas
automatiquement ce stash sur la branche suivante.

\subsection{4. Exercice 1 --- Enquête : contrôle d'accès cassé (10
min)}\label{exercice-1-enquuxeate-contruxf4le-daccuxe8s-cassuxe9-10-min}

\textbf{Branche :} \texttt{Tp\_version1}. Dans cette première étape, ne
corrigez pas le code : votre objectif est d'observer, de démontrer et
d'expliquer.

\subsubsection{Point de départ : OWASP Top
10}\label{point-de-duxe9part-owasp-top-10}

L'\href{https://top10.owasp.org/2025/A01_2025-Broken_Access_Control/}{OWASP
Top 10:2025, A01 --- Broken Access Control} décrit les situations où une
application laisse quelqu'un agir au-delà des permissions prévues. Parmi
les exemples cités : une API qui n'applique pas de contrôle aux
opérations \texttt{POST}, \texttt{PUT} ou \texttt{DELETE}, ou un accès à
une fonction privilégiée en appelant directement son URL. Dans cet
exercice, recherchez des indices de cette catégorie dans Planora. Vous
devrez établir vous-même quels accès sont problématiques et proposer les
règles à appliquer.

\subsubsection{Mener l'enquête}\label{mener-lenquuxeate}

\begin{enumerate}
\def\labelenumi{\arabic{enumi}.}
\tightlist
\item
  Lancez Spring Boot et Planora. Restez \textbf{déconnecté}.
\item
  Dans \textbf{Mes tâches}, consultez les tâches, puis créez-en une,
  modifiez-la, terminez-la et supprimez-la. Notez ce que l'interface
  vous laisse faire.
\item
  Ouvrez \textbf{Statistiques}. La navigation dans React est un indice à
  relever, pas une preuve suffisante du contrôle d'accès.
\item
  Vérifiez les requêtes directement contre l'API : observez les appels
  dans l'onglet \textbf{Réseau} du navigateur ou utilisez \texttt{curl}.
  Testez les routes listées dans la fiche ci-dessous \textbf{sans
  identifiants}. Pour les opérations de modification, utilisez un
  identifiant de tâche existant ; supprimez uniquement la tâche que vous
  venez de créer.
\item
  Comparez ce que vous observez avec les permissions que vous estimez
  nécessaires pour une application d'équipe. Pour vous guider, regardez
  qui peut lire les tâches, modifier les données de l'équipe et
  consulter des statistiques d'administration.
\item
  Lisez \texttt{TaskController}, \texttt{AdminController},
  \texttt{TaskService} et \texttt{StatisticsService}. Repérez où sont
  déclarées les routes et où se trouve (ou ne se trouve pas) une
  vérification d'accès.
\end{enumerate}

L'interface publique \texttt{/api/public/info} peut être vérifiée avec
\texttt{curl} ; elle n'a pas d'écran dédié dans Planora. Les réponses
initiales de l'API sont observables ainsi :

\begin{Shaded}
\begin{Highlighting}[]
\ExtensionTok{curl} \AttributeTok{{-}i}\NormalTok{ http://localhost:8080/api/tasks}
\ExtensionTok{curl} \AttributeTok{{-}i}\NormalTok{ http://localhost:8080/api/admin/stats}
\ExtensionTok{curl} \AttributeTok{{-}i}\NormalTok{ http://localhost:8080/api/public/info}
\end{Highlighting}
\end{Shaded}

Ne vous contentez pas de dire qu'une opération « fonctionne » : gardez
le statut HTTP, le résultat visible et l'action réalisée comme éléments
de preuve.

\subsubsection{Fiche d'observation à
compléter}\label{fiche-dobservation-uxe0-compluxe9ter}

La fiche ci-dessous fait partie de ce support. Complétez-la dans le PDF
(ou imprimez ces pages). Pour chaque requête, notez ce que \textbf{vous}
avez observé sans vous connecter. Si une réponse diffère de ce que vous
attendiez, expliquez pourquoi cela vous semble important.

{\def\LTcaptype{none} % do not increment counter
\begin{longtable}[]{@{}
  >{\raggedright\arraybackslash}p{(\linewidth - 6\tabcolsep) * \real{0.2500}}
  >{\raggedright\arraybackslash}p{(\linewidth - 6\tabcolsep) * \real{0.2500}}
  >{\raggedright\arraybackslash}p{(\linewidth - 6\tabcolsep) * \real{0.2500}}
  >{\raggedright\arraybackslash}p{(\linewidth - 6\tabcolsep) * \real{0.2500}}@{}}
\toprule\noalign{}
\begin{minipage}[b]{\linewidth}\raggedright
Requête testée sans connexion
\end{minipage} & \begin{minipage}[b]{\linewidth}\raggedright
Accessible ?
\end{minipage} & \begin{minipage}[b]{\linewidth}\raggedright
Statut HTTP
\end{minipage} & \begin{minipage}[b]{\linewidth}\raggedright
Résultat ou effet observé
\end{minipage} \\
\midrule\noalign{}
\endhead
\bottomrule\noalign{}
\endlastfoot
\texttt{GET\ /api/tasks} & \_\_\_\_\_\_\_\_\_\_ & \_\_\_\_\_\_ &
\_\_\_\_\_\_\_\_\_\_\_\_\_\_\_\_\_\_\_\_\_\_\_\_\_\_\_\_\_\_\_\_\_\_\_\_\_\_\_\_\_\_ \\
\texttt{POST\ /api/tasks} & \_\_\_\_\_\_\_\_\_\_ & \_\_\_\_\_\_ &
\_\_\_\_\_\_\_\_\_\_\_\_\_\_\_\_\_\_\_\_\_\_\_\_\_\_\_\_\_\_\_\_\_\_\_\_\_\_\_\_\_\_ \\
\texttt{PUT\ /api/tasks/\{id\}} & \_\_\_\_\_\_\_\_\_\_ & \_\_\_\_\_\_ &
\_\_\_\_\_\_\_\_\_\_\_\_\_\_\_\_\_\_\_\_\_\_\_\_\_\_\_\_\_\_\_\_\_\_\_\_\_\_\_\_\_\_ \\
\texttt{PATCH\ /api/tasks/\{id\}/complete} & \_\_\_\_\_\_\_\_\_\_ &
\_\_\_\_\_\_ &
\_\_\_\_\_\_\_\_\_\_\_\_\_\_\_\_\_\_\_\_\_\_\_\_\_\_\_\_\_\_\_\_\_\_\_\_\_\_\_\_\_\_ \\
\texttt{DELETE\ /api/tasks/\{id\}} & \_\_\_\_\_\_\_\_\_\_ & \_\_\_\_\_\_
&
\_\_\_\_\_\_\_\_\_\_\_\_\_\_\_\_\_\_\_\_\_\_\_\_\_\_\_\_\_\_\_\_\_\_\_\_\_\_\_\_\_\_ \\
\texttt{GET\ /api/admin/stats} & \_\_\_\_\_\_\_\_\_\_ & \_\_\_\_\_\_ &
\_\_\_\_\_\_\_\_\_\_\_\_\_\_\_\_\_\_\_\_\_\_\_\_\_\_\_\_\_\_\_\_\_\_\_\_\_\_\_\_\_\_ \\
\end{longtable}
}

\textbf{1. Quels droits proposeriez-vous ?} Indiquez ce qu'une personne
sans compte, un membre connecté et un administrateur devraient pouvoir
faire dans cette application. Justifiez brièvement.

\begin{itemize}
\tightlist
\item
  Une personne sans compte :
  \_\_\_\_\_\_\_\_\_\_\_\_\_\_\_\_\_\_\_\_\_\_\_\_\_\_\_\_\_\_\_\_\_\_\_\_\_\_\_\_\_\_\_\_\_\_\_\_\_\_\_\_
\item
  Un membre connecté de l'équipe :
  \_\_\_\_\_\_\_\_\_\_\_\_\_\_\_\_\_\_\_\_\_\_\_\_\_\_\_\_\_\_\_\_\_\_\_\_\_\_\_\_\_\_\_\_\_\_
\item
  Un administrateur :
  \_\_\_\_\_\_\_\_\_\_\_\_\_\_\_\_\_\_\_\_\_\_\_\_\_\_\_\_\_\_\_\_\_\_\_\_\_\_\_\_\_\_\_\_\_\_\_\_\_\_\_\_\_\_\_\_\_\_\_
\end{itemize}

\textbf{2. Que vous apprend l'enquête sur OWASP A01 ?} Comparez les
droits proposés aux résultats observés. Donnez deux exemples précis de
requêtes qui vous semblent problématiques et expliquez le risque
associé.

\begin{center}\rule{0.5\linewidth}{0.5pt}\end{center}

\begin{center}\rule{0.5\linewidth}{0.5pt}\end{center}

\begin{center}\rule{0.5\linewidth}{0.5pt}\end{center}

\textbf{3. Où faut-il vérifier les droits ?} Expliquez avec vos mots la
différence entre authentification et autorisation. Pourquoi le contrôle
doit-il être appliqué dans le backend, même si React peut masquer des
boutons ?

\begin{center}\rule{0.5\linewidth}{0.5pt}\end{center}

\begin{center}\rule{0.5\linewidth}{0.5pt}\end{center}

\begin{center}\rule{0.5\linewidth}{0.5pt}\end{center}

\textbf{Conclusion ---} En une phrase, reliez vos observations à
\textbf{A01:2025 Broken Access Control}. Gardez cette fiche pour
comparer les résultats après les exercices suivants.

\subsection{5. Exercice 2 --- Activer la protection par défaut (10
min)}\label{exercice-2-activer-la-protection-par-duxe9faut-10-min}

\textbf{Branche :} \texttt{Tp\_version2}. \textbf{Fichier à modifier :}
\texttt{pom.xml} uniquement.

\subsubsection{Objectif}\label{objectif}

Ajouter Spring Security au projet et observer ce que Spring Boot met en
place automatiquement, sans écrire encore de règles personnalisées. Vous
allez constater que la protection par défaut demande une
authentification pour toutes les routes, même celle nommée
\texttt{/api/public/info}. Spring crée aussi temporairement le compte
\texttt{user} et affiche son mot de passe dans la console.

À la fin de l'exercice, vous devriez pouvoir expliquer ce que la
dépendance change, comment retrouver le compte de développement et
pourquoi les rôles (USER/ADMIN) ne sont pas encore distingués.

\subsubsection{TODO 2.1 --- Ajouter la
dépendance}\label{todo-2.1-ajouter-la-duxe9pendance}

Dans \texttt{pom.xml}, repérez le commentaire
\texttt{TODO\ EXERCICE\ 2}. Remplacez-le par ce bloc Maven :

\begin{Shaded}
\begin{Highlighting}[]
\NormalTok{\textless{}}\KeywordTok{dependency}\NormalTok{\textgreater{}}
\NormalTok{    \textless{}}\KeywordTok{groupId}\NormalTok{\textgreater{}org.springframework.boot\textless{}/}\KeywordTok{groupId}\NormalTok{\textgreater{}}
\NormalTok{    \textless{}}\KeywordTok{artifactId}\NormalTok{\textgreater{}spring{-}boot{-}starter{-}security\textless{}/}\KeywordTok{artifactId}\NormalTok{\textgreater{}}
\NormalTok{\textless{}/}\KeywordTok{dependency}\NormalTok{\textgreater{}}
\end{Highlighting}
\end{Shaded}

N'ajoutez pas de balise \texttt{\textless{}version\textgreater{}} : le
parent Spring Boot gère cette version. N'ajoutez pas de configuration
Java et ne modifiez pas les contrôleurs. À cette étape, il suffit
d'ajouter le starter.

\subsubsection{TODO 2.2 --- Observer l'API sans
identifiants}\label{todo-2.2-observer-lapi-sans-identifiants}

Redémarrez Spring avec \texttt{mvn\ spring-boot:run}. Dans Planora, sans
vous connecter, consultez \textbf{Mes tâches} puis
\textbf{Statistiques}. Relevez les réponses dans l'onglet
\textbf{Réseau} du navigateur. L'information \texttt{/api/public/info}
n'a pas d'écran dans React ; vérifiez-la directement :

\begin{Shaded}
\begin{Highlighting}[]
\ExtensionTok{curl} \AttributeTok{{-}i}\NormalTok{ http://localhost:8080/api/public/info}
\ExtensionTok{curl} \AttributeTok{{-}i}\NormalTok{ http://localhost:8080/api/tasks}
\ExtensionTok{curl} \AttributeTok{{-}i}\NormalTok{ http://localhost:8080/api/admin/stats}
\end{Highlighting}
\end{Shaded}

\textbf{À observer :} les requêtes de lecture renvoient
\texttt{401\ Unauthorized}, y compris l'information publique. Expliquez
pourquoi un nom de route contenant \texttt{public} ne suffit pas à la
rendre accessible.

\subsubsection{TODO 2.3 --- Retrouver l'identité
temporaire}\label{todo-2.3-retrouver-lidentituxe9-temporaire}

Dans la console du backend, repérez le message Spring qui affiche le
\textbf{mot de passe généré}. Il change au démarrage : utilisez celui de
l'exécution en cours et ne le copiez pas dans le dépôt.

Dans Planora, ouvrez \textbf{Connexion}, saisissez l'identifiant
\texttt{user} et ce mot de passe, puis consultez les tâches. Vérifiez
aussi la requête avec curl :

\begin{Shaded}
\begin{Highlighting}[]
\ExtensionTok{curl} \AttributeTok{{-}i} \AttributeTok{{-}u}\NormalTok{ user:MOT\_DE\_PASSE\_GENERE http://localhost:8080/api/tasks}
\end{Highlighting}
\end{Shaded}

Remplacez \texttt{MOT\_DE\_PASSE\_GENERE} par la valeur de votre
console. La lecture doit fonctionner après authentification. Le fait que
la page React affiche \texttt{user} ne suffit pas à prouver la connexion
: vérifiez la réponse de l'API.

\subsubsection{TODO 2.4 --- Comparer une lecture et une
mutation}\label{todo-2.4-comparer-une-lecture-et-une-mutation}

Essayez de créer une tâche depuis Planora avec le compte \texttt{user},
puis notez le statut HTTP. Spring Security active aussi la protection
CSRF par défaut ; le client pédagogique n'envoie pas de jeton CSRF. Une
mutation peut donc recevoir \texttt{403} même avec le bon mot de passe.
Les lectures \texttt{GET}, elles, sont acceptées avec les identifiants
valides.

Ne désactivez pas CSRF dans cet exercice : cette décision fait partie de
la configuration de l'API à l'exercice 4. Pour l'instant, distinguez le
refus d'authentification (\texttt{401}) du refus d'une mutation sans
jeton CSRF (\texttt{403}).

\subsubsection{TODO 2.5 --- Vérifier avec le test
fourni}\label{todo-2.5-vuxe9rifier-avec-le-test-fourni}

Dans
\texttt{src/test/java/fr/tpspringsecurity/TaskApiSecurityTest.java},
repérez \texttt{exercice2LesRoutesDemandentUneConnexion}. Après avoir
ajouté la dépendance, retirez l'annotation \texttt{@Disabled} juste
au-dessus de cette méthode, puis exécutez :

\begin{Shaded}
\begin{Highlighting}[]
\ExtensionTok{mvn}\NormalTok{ test}
\end{Highlighting}
\end{Shaded}

Le test vérifie que les trois lectures (\texttt{/api/public/info},
\texttt{/api/tasks}, \texttt{/api/admin/stats}) refusent un visiteur
anonyme. Il ne vérifie pas les mutations.

\subsubsection{Questions de
compréhension}\label{questions-de-compruxe9hension}

\begin{enumerate}
\def\labelenumi{\arabic{enumi}.}
\tightlist
\item
  Qui crée le compte \texttt{user}, puisque vous ne l'avez pas déclaré
  dans le code ?
\item
  À cette étape, quelle règle permet de distinguer un utilisateur d'un
  administrateur ? Appuyez-vous sur ce que vous avez observé.
\item
  Pourquoi la route d'information publique exige-t-elle aussi une
  authentification ?
\item
  Pourquoi une mutation peut-elle renvoyer \texttt{403} alors que votre
  lecture authentifiée renvoie \texttt{200} ?
\end{enumerate}

\textbf{Pour continuer :} la branche \texttt{Tp\_version3} garde la
protection par défaut et remplace le compte temporaire par Alice et Bob.

\subsection{6. Exercice 3 --- Déclarer les comptes avec BCrypt (15
min)}\label{exercice-3-duxe9clarer-les-comptes-avec-bcrypt-15-min}

\textbf{Branche :} \texttt{Tp\_version3}. \textbf{Fichier à modifier :}
\texttt{src/main/java/fr/tpspringsecurity/config/UserConfig.java}.

\subsubsection{Objectif}\label{objectif-1}

Remplacer le compte temporaire généré par Spring par deux comptes de TP
définis dans l'application, puis stocker leurs mots de passe sous forme
BCrypt. Vous configurez \textbf{qui peut s'authentifier} et les rôles de
ces comptes. Vous ne configurez pas encore \textbf{quelles routes chaque
rôle peut consulter} : c'est l'exercice 4.

{\def\LTcaptype{none} % do not increment counter
\begin{longtable}[]{@{}lll@{}}
\toprule\noalign{}
Identifiant exact & Mot de passe de démonstration & Rôle à déclarer \\
\midrule\noalign{}
\endhead
\bottomrule\noalign{}
\endlastfoot
\texttt{alice} & \texttt{alice123} & \texttt{USER} \\
\texttt{bob} & \texttt{bob123} & \texttt{ADMIN} \\
\end{longtable}
}

\subsubsection{Avant de coder}\label{avant-de-coder}

Dans cette branche, Spring Security est déjà installé et sa politique
HTTP par défaut est active. Au début, les routes demandent un compte
\texttt{user} généré. \texttt{UserConfig.java} est volontairement
incomplet ; ses commentaires indiquent les TODO 3.1 à 3.4. Travaillez
uniquement dans ce fichier. Ne créez pas de \texttt{SecurityFilterChain}
et ne modifiez pas \texttt{SecurityConfig.java}.

\subsubsection{TODO 3.1 --- Déclarer l'encodeur
BCrypt}\label{todo-3.1-duxe9clarer-lencodeur-bcrypt}

Créez un Bean de type \texttt{PasswordEncoder}. Il doit retourner une
instance de \texttt{BCryptPasswordEncoder}. Cet objet servira à encoder
les mots de passe au moment de créer les comptes, puis Spring s'en
servira pour vérifier ceux saisis à la connexion.

Imports à repérer : \texttt{Bean}, \texttt{PasswordEncoder} et
\texttt{BCryptPasswordEncoder}.

\subsubsection{TODO 3.2 --- Déclarer la source des
utilisateurs}\label{todo-3.2-duxe9clarer-la-source-des-utilisateurs}

Créez un Bean \texttt{UserDetailsService} qui reçoit le
\texttt{PasswordEncoder} en paramètre :

\begin{Shaded}
\begin{Highlighting}[]
\AttributeTok{@Bean}
\KeywordTok{public}\NormalTok{ UserDetailsService }\FunctionTok{userDetailsService}\OperatorTok{(}\NormalTok{PasswordEncoder encoder}\OperatorTok{)} \OperatorTok{\{}
    \CommentTok{// Créer les deux comptes puis les stocker en mémoire.}
\OperatorTok{\}}
\end{Highlighting}
\end{Shaded}

Spring injectera le Bean créé au TODO 3.1. La méthode doit fournir au
système de sécurité les comptes qu'il utilisera pour authentifier les
requêtes.

\subsubsection{TODO 3.3 --- Construire les deux
comptes}\label{todo-3.3-construire-les-deux-comptes}

Dans cette méthode, construisez un \texttt{UserDetails} pour Alice et un
pour Bob avec \texttt{User.withUsername(...)}. Pour chacun, renseignez
l'identifiant du tableau, un mot de passe encodé et le rôle
correspondant (\texttt{USER} ou \texttt{ADMIN}). Terminez chaque
construction avec \texttt{.build()}.

\textbf{Attention :} passez le mot de passe à
\texttt{encoder.encode(...)} avant de l'attribuer. Ne stockez pas le
texte \texttt{alice123} ou \texttt{bob123} directement dans l'objet
utilisateur. Utilisez \texttt{.roles("ADMIN")} pour déclarer le rôle de
Bob ; Spring le convertit en autorité \texttt{ROLE\_ADMIN}.

Imports à repérer : \texttt{User}, \texttt{UserDetails} et
\texttt{UserDetailsService}.

\subsubsection{TODO 3.4 --- Retourner le gestionnaire en
mémoire}\label{todo-3.4-retourner-le-gestionnaire-en-muxe9moire}

Retournez un \texttt{InMemoryUserDetailsManager} contenant les deux
objets construits. Il remplace la source d'utilisateurs automatique pour
l'application : après redémarrage, le compte \texttt{user} généré n'est
plus celui que vous utilisez.

Import à repérer : \texttt{InMemoryUserDetailsManager}.

\subsubsection{Vérifier dans Planora et avec
l'API}\label{vuxe9rifier-dans-planora-et-avec-lapi}

Redémarrez Spring. Essayez d'abord l'ancien compte \texttt{user}, puis
connectez-vous avec Alice et Bob. Pour chacun, consultez \textbf{Mes
tâches} et \textbf{Statistiques}. Essayez enfin Alice avec un mauvais
mot de passe. Vous pouvez contrôler une lecture avec curl :

\begin{Shaded}
\begin{Highlighting}[]
\ExtensionTok{curl} \AttributeTok{{-}i} \AttributeTok{{-}u}\NormalTok{ alice:alice123 http://localhost:8080/api/tasks}
\ExtensionTok{curl} \AttributeTok{{-}i} \AttributeTok{{-}u}\NormalTok{ bob:bob123 http://localhost:8080/api/admin/stats}
\ExtensionTok{curl} \AttributeTok{{-}i} \AttributeTok{{-}u}\NormalTok{ alice:incorrect http://localhost:8080/api/tasks}
\end{Highlighting}
\end{Shaded}

\textbf{Résultats à observer :} Alice et Bob obtiennent \texttt{200} sur
les tâches et les statistiques ; le mauvais mot de passe reçoit
\texttt{401}. Les statistiques répondent encore \texttt{200} pour Alice
: son rôle est bien enregistré, mais aucune règle HTTP ne l'utilise
encore. Le chemin \texttt{/admin/} ne réserve pas une route à un
administrateur à lui seul.

Les mutations peuvent répondre \texttt{403} même quand Alice ou Bob est
authentifié. C'est la protection CSRF par défaut de cette étape, pas une
conséquence de leur rôle. Ne changez pas CSRF ici ; l'exercice 4
configurera les règles HTTP et le comportement des mutations.

\subsubsection{TODO 3.5 --- Activer le test
fourni}\label{todo-3.5-activer-le-test-fourni}

Dans
\texttt{src/test/java/fr/tpspringsecurity/TaskApiSecurityTest.java},
repérez \texttt{exercice3AliceEtBobPeuventSeConnecter}. Retirez
\texttt{@Disabled} située juste au-dessus de cette méthode et exécutez :

\begin{Shaded}
\begin{Highlighting}[]
\ExtensionTok{mvn}\NormalTok{ test}
\end{Highlighting}
\end{Shaded}

Ce test vérifie qu'Alice et Bob lisent la liste des tâches
(\texttt{200}) et qu'un mauvais mot de passe pour Alice est refusé
(\texttt{401}). Il ne vérifie pas encore que les statistiques sont
réservées à Bob : ajoutez cette observation depuis Planora ou curl.

\subsubsection{Questions de
compréhension}\label{questions-de-compruxe9hension-1}

\begin{enumerate}
\def\labelenumi{\arabic{enumi}.}
\tightlist
\item
  Quel rôle joue \texttt{PasswordEncoder} à la création du compte puis
  lors d'une connexion ?
\item
  Le rôle \texttt{ADMIN} protège-t-il déjà \texttt{/api/admin/stats} ?
  Justifiez avec votre observation.
\item
  Quelle différence faites-vous entre authentifier un compte et
  autoriser ce compte à utiliser une route ?
\end{enumerate}

\textbf{Pour continuer :} la branche \texttt{Tp\_version4} contient la
correction de \texttt{UserConfig.java}. Vous y configurerez les
autorisations HTTP.

\subsection{7. Exercice 4 --- Définir les autorisations HTTP (10
min)}\label{exercice-4-duxe9finir-les-autorisations-http-10-min}

\textbf{Branche :} \texttt{Tp\_version4}. \textbf{Fichier à modifier :}
\texttt{src/main/java/fr/tpspringsecurity/config/SecurityConfig.java}.

\subsubsection{Objectif}\label{objectif-2}

Les comptes Alice (\texttt{USER}) et Bob (\texttt{ADMIN}) existent et
peuvent tous deux lire les tâches et les statistiques. Vous allez
maintenant faire appliquer leurs permissions par Spring Security :
laisser une route explicitement publique, limiter les statistiques aux
administrateurs, puis exiger une authentification pour le reste de
l'API.

Vous allez également configurer HTTP Basic sans session, et adapter CSRF
afin que les mutations du client React fourni fonctionnent. Les comptes
et leurs rôles restent dans \texttt{UserConfig.java}.

\subsubsection{Avant de coder}\label{avant-de-coder-1}

Dans \texttt{SecurityConfig.java}, les TODO 4.1 à 4.6 sont des étapes
d'une seule chaîne de filtres. Les imports nécessaires ne sont pas tous
présents : ajoutez-les depuis les types indiqués ci-dessous. Gardez les
règles d'accès dans le même ordre que les étapes.

\subsubsection{TODO 4.1 --- Déclarer la chaîne de
sécurité}\label{todo-4.1-duxe9clarer-la-chauxeene-de-suxe9curituxe9}

Ajoutez l'annotation \texttt{@Bean} sur une méthode qui retourne
\texttt{SecurityFilterChain} et reçoit \texttt{HttpSecurity}. La méthode
doit configurer puis retourner la chaîne ; c'est là que Spring décide si
une requête anonyme ou authentifiée peut atteindre les contrôleurs.

Types à importer : \texttt{Bean}, \texttt{HttpSecurity},
\texttt{SecurityFilterChain}.

\subsubsection{TODO 4.2 à 4.4 --- Définir et ordonner les règles
d'accès}\label{todo-4.2-uxe0-4.4-duxe9finir-et-ordonner-les-ruxe8gles-daccuxe8s}

Dans \texttt{authorizeHttpRequests}, ajoutez les règles dans cet ordre :

\begin{enumerate}
\def\labelenumi{\arabic{enumi}.}
\tightlist
\item
  Pour \texttt{GET\ /api/public/info}, utilisez
  \texttt{requestMatchers(HttpMethod.GET,\ "/api/public/info").permitAll()}.
\item
  Pour \texttt{GET\ /api/admin/**}, utilisez
  \texttt{requestMatchers(HttpMethod.GET,\ "/api/admin/**").hasRole("ADMIN")}.
\item
  Pour toutes les requêtes restantes, utilisez
  \texttt{anyRequest().authenticated()}.
\end{enumerate}

Importez \texttt{HttpMethod}. Les règles précises doivent être évaluées
avant \texttt{anyRequest()}, qui termine la liste.
\texttt{hasRole("ADMIN")} contrôle l'autorité \texttt{ROLE\_ADMIN},
ajoutée aux comptes avec \texttt{.roles("ADMIN")} à l'exercice
précédent.

\subsubsection{TODO 4.5 --- Configurer HTTP Basic et les
sessions}\label{todo-4.5-configurer-http-basic-et-les-sessions}

Configurez la politique de session avec
\texttt{SessionCreationPolicy.STATELESS}, puis activez HTTP Basic avec
\texttt{httpBasic(withDefaults())}. Chaque requête protégée doit alors
présenter ses identifiants ; le serveur ne conserve pas une session de
connexion entre les appels.

Types à importer : \texttt{SessionCreationPolicy} et
\texttt{withDefaults} de
\texttt{org.springframework.security.config.Customizer}.

\subsubsection{TODO 4.6 --- Adapter CSRF au client
fourni}\label{todo-4.6-adapter-csrf-au-client-fourni}

Ajoutez \texttt{.csrf(csrf\ -\textgreater{}\ csrf.disable())} dans la
chaîne avant de la construire avec \texttt{.build()}.

Ce choix permet au client de TP d'envoyer ses requêtes JSON POST, PUT,
PATCH et DELETE avec l'en-tête \texttt{Authorization} qu'il construit
explicitement, et sans cookie. \textbf{Ne généralisez pas cette
désactivation à toutes les API stateless ou HTTP Basic} : le navigateur
peut envoyer automatiquement des identifiants Basic dans d'autres
configurations, ce qui peut exposer à CSRF. Si le mode
d'authentification ou le client change, réévaluez la protection.
Consultez la
\href{https://docs.spring.io/spring-security/reference/features/exploits/csrf.html}{documentation
Spring Security sur CSRF}.

\subsubsection{Vérifier les permissions dans
Planora}\label{vuxe9rifier-les-permissions-dans-planora}

Redémarrez Spring puis vérifiez les cas suivants. Relevez les statuts
dans l'onglet Réseau ; les opérations métier doivent réellement réussir
ou être refusées par le serveur.

\begin{itemize}
\tightlist
\item
  \textbf{Sans connexion :} Mes tâches est refusé avec \texttt{401}. La
  route publique \texttt{/api/public/info} doit être vérifiée avec curl
  et renvoyer \texttt{200}.
\item
  \textbf{Alice :} elle peut consulter les tâches, créer (\texttt{201}),
  modifier et terminer (\texttt{200}), puis supprimer (\texttt{204}).
  L'ouverture des statistiques est refusée avec \texttt{403}.
\item
  \textbf{Bob :} il peut gérer les tâches et consulter les statistiques
  (\texttt{200}).
\item
  \textbf{Projet confidentiel :} dans cette branche, la route n'existe
  pas encore. Sans compte, Spring refuse d'abord la requête
  (\texttt{401}). Avec Alice ou Bob, la requête passe le filtre HTTP,
  puis le contrôleur absent entraîne \texttt{404}.
\end{itemize}

\begin{Shaded}
\begin{Highlighting}[]
\ExtensionTok{curl} \AttributeTok{{-}i}\NormalTok{ http://localhost:8080/api/public/info}
\ExtensionTok{curl} \AttributeTok{{-}i} \AttributeTok{{-}u}\NormalTok{ alice:alice123 http://localhost:8080/api/admin/stats}
\ExtensionTok{curl} \AttributeTok{{-}i} \AttributeTok{{-}u}\NormalTok{ bob:bob123 http://localhost:8080/api/admin/stats}
\end{Highlighting}
\end{Shaded}

\subsubsection{Activer le test fourni}\label{activer-le-test-fourni}

Dans \texttt{TaskApiSecurityTest.java}, retirez \texttt{@Disabled}
au-dessus de
\texttt{exercice4LAccesPublicEtLesDroitsAdminSontRespectes}, puis lancez
\texttt{mvn\ test}. Le test vérifie la route publique, le refus des
statistiques à Alice, leur accès par Bob et la création d'une tâche par
Alice. Cette dernière requête permet aussi de vérifier que CSRF a bien
été configuré pour le client.

Les autres tests actifs vérifient déjà qu'un visiteur ne lit pas les
tâches (\texttt{401}) et que les deux comptes valides y accèdent
(\texttt{200}).

\subsubsection{Questions de
compréhension}\label{questions-de-compruxe9hension-2}

\begin{enumerate}
\def\labelenumi{\arabic{enumi}.}
\tightlist
\item
  Pourquoi faut-il placer les règles publiques et ADMIN avant
  \texttt{anyRequest()} ?
\item
  Pourquoi Alice reçoit-elle \texttt{403} sur les statistiques, alors
  qu'elle est authentifiée ?
\item
  Que changent \texttt{STATELESS} et la désactivation de CSRF pour les
  requêtes de ce client ? Quelle précaution garder en tête si le client
  évolue ?
\end{enumerate}

\textbf{Pour continuer :} la branche \texttt{Tp\_version5} contient la
chaîne HTTP complète. Vous y protégerez une méthode Java avec
\texttt{@PreAuthorize}.

\subsection{\texorpdfstring{8. Exercice 5 --- Protéger une méthode avec
\texttt{@PreAuthorize} (10
min)}{8. Exercice 5 --- Protéger une méthode avec @PreAuthorize (10 min)}}\label{exercice-5-protuxe9ger-une-muxe9thode-avec-preauthorize-10-min}

\textbf{Branche :} \texttt{Tp\_version5}. \textbf{Fichiers à modifier :}
\texttt{src/main/java/fr/tpspringsecurity/controller/TaskController.java}
et
\texttt{src/main/java/fr/tpspringsecurity/config/SecurityConfig.java}.

\subsubsection{Objectif}\label{objectif-3}

La chaîne HTTP de l'exercice 4 protège déjà les statistiques, mais elle
exige seulement une authentification pour les routes restantes sous
\texttt{/api/tasks}. Vous allez ajouter une deuxième vérification, au
niveau d'une méthode Java, pour que seul le rôle \texttt{ADMIN} puisse
lire le projet confidentiel.

Cet exercice montre la différence entre filtrer une requête HTTP selon
son chemin et vérifier les permissions juste avant l'exécution d'une
méthode. La ressource métier et sa route sont déjà présentes : il ne
faut pas les recréer.

\subsubsection{État initial à
constater}\label{uxe9tat-initial-uxe0-constater}

La branche 5 expose déjà \texttt{GET\ /api/tasks/confidential}. Le
contrôleur appelle \texttt{taskService.confidential()} et retourne un
objet \texttt{ConfidentialProject}. À l'ouverture de la branche, la
règle HTTP générale demande une authentification, mais Alice peut encore
lire la note (\texttt{200}). Bob y accède également (\texttt{200}) ; un
visiteur reçoit \texttt{401}.

Connectez-vous comme Alice dans Planora et ouvrez \textbf{Projet
confidentiel} avant la correction. Relevez le statut. Cette requête est
votre état de comparaison.

\subsubsection{TODO 5.1 --- Protéger la méthode du
contrôleur}\label{todo-5.1-protuxe9ger-la-muxe9thode-du-contruxf4leur}

Dans \texttt{TaskController.java}, repérez la méthode existante annotée
\texttt{@GetMapping("/confidential")} et le commentaire juste au-dessus.

\begin{itemize}
\tightlist
\item
  Importez
  \texttt{org.springframework.security.access.prepost.PreAuthorize} si
  nécessaire.
\item
  Activez l'annotation
  \texttt{@PreAuthorize("hasRole(\textquotesingle{}ADMIN\textquotesingle{})")}
  sur la méthode \texttt{confidential()}.
\item
  Laissez son corps appeler le service comme il le fait déjà.
\end{itemize}

Ne créez pas un nouvel endpoint et ne déplacez pas la logique métier
dans le contrôleur.

\subsubsection{TODO 5.2 --- Activer la sécurité des
méthodes}\label{todo-5.2-activer-la-suxe9curituxe9-des-muxe9thodes}

Dans \texttt{SecurityConfig.java}, ajoutez
\texttt{@EnableMethodSecurity} à la classe de configuration. Sans cette
annotation, la vérification \texttt{@PreAuthorize} n'est pas activée.

N'ajoutez pas de matcher HTTP pour \texttt{/api/tasks/confidential}. La
règle générale existante exige déjà une identité ; dans cet exercice, la
restriction supplémentaire au rôle ADMIN doit venir de la méthode.

\subsubsection{Vérifier après
redémarrage}\label{vuxe9rifier-apruxe8s-reduxe9marrage}

Réessayez la route dans Planora avec chaque identité et observez les
réponses :

{\def\LTcaptype{none} % do not increment counter
\begin{longtable}[]{@{}ll@{}}
\toprule\noalign{}
Profil & Résultat attendu \\
\midrule\noalign{}
\endhead
\bottomrule\noalign{}
\endlastfoot
Sans compte & \texttt{401} --- authentification exigée par la règle
HTTP \\
Alice (\texttt{USER}) & \texttt{403} --- identifiée, mais le rôle ne
suffit pas \\
Bob (\texttt{ADMIN}) & \texttt{200} --- méthode autorisée \\
\end{longtable}
}

\texttt{hasRole("ADMIN")} recherche l'autorité \texttt{ROLE\_ADMIN}, que
Spring associe au rôle déclaré avec \texttt{.roles("ADMIN")}.

\subsubsection{TODO 5.3 --- Activer le test
fourni}\label{todo-5.3-activer-le-test-fourni}

Dans \texttt{TaskApiSecurityTest.java}, retirez \texttt{@Disabled} juste
au-dessus de \texttt{exercice5MethodeConfidentielleReserveeAuxAdmins},
puis lancez \texttt{mvn\ test}. Le test vérifie les trois cas de la
matrice (anonyme, Alice, Bob).

Vérifiez aussi que Mes tâches et Statistiques conservent les droits de
l'exercice précédent : Alice peut gérer les tâches mais reçoit
\texttt{403} sur les statistiques ; Bob peut accéder aux deux.

\subsubsection{Questions de
compréhension}\label{questions-de-compruxe9hension-3}

\begin{enumerate}
\def\labelenumi{\arabic{enumi}.}
\tightlist
\item
  Quel est l'état de l'accès d'Alice avant et après l'ajout de
  \texttt{@PreAuthorize} ?
\item
  Que fait \texttt{@EnableMethodSecurity} ?
\item
  Quelle différence entre la règle générale \texttt{authenticated()} et
  la vérification de rôle sur cette méthode ? Pourquoi les données
  restent-elles dans le service ?
\end{enumerate}

\textbf{Pour continuer :} la branche \texttt{Tp\_version6} contient les
deux annotations actives et les tests complets.

\subsection{9. Exercice 6 --- Vérifier le parcours final (5
min)}\label{exercice-6-vuxe9rifier-le-parcours-final-5-min}

\textbf{Branche :} \texttt{Tp\_version6}. La correction complète est
fournie ; aucun TODO de sécurité à compléter.

\begin{Shaded}
\begin{Highlighting}[]
\ExtensionTok{mvn}\NormalTok{ clean verify}
\end{Highlighting}
\end{Shaded}

Le test d'intégration couvre l'accès public, les identifiants invalides,
les accès aux statistiques et à la note, les mutations d'Alice,
l'évolution des statistiques et la suppression d'une tâche inexistante.
Les tests unitaires vérifient la logique métier. Tous les tests sont
actifs en version 6.

La matrice suivante suppose des données valides et un identifiant de
tâche existant :

{\def\LTcaptype{none} % do not increment counter
\begin{longtable}[]{@{}lrrr@{}}
\toprule\noalign{}
Opération & Anonyme & Alice \texttt{USER} & Bob \texttt{ADMIN} \\
\midrule\noalign{}
\endhead
\bottomrule\noalign{}
\endlastfoot
\texttt{GET\ /api/public/info} & \texttt{200} & \texttt{200} &
\texttt{200} \\
\texttt{GET\ /api/tasks} & \texttt{401} & \texttt{200} & \texttt{200} \\
\texttt{POST\ /api/tasks} & \texttt{401} & \texttt{201} &
\texttt{201} \\
\texttt{PUT\ /api/tasks/\{id\}} & \texttt{401} & \texttt{200} &
\texttt{200} \\
\texttt{PATCH\ /api/tasks/\{id\}/complete} & \texttt{401} & \texttt{200}
& \texttt{200} \\
\texttt{DELETE\ /api/tasks/\{id\}} & \texttt{401} & \texttt{204} &
\texttt{204} \\
\texttt{GET\ /api/admin/stats} & \texttt{401} & \texttt{403} &
\texttt{200} \\
\texttt{GET\ /api/tasks/confidential} & \texttt{401} & \texttt{403} &
\texttt{200} \\
\end{longtable}
}

Dans Planora, refaites un cycle complet de gestion des tâches avec
Alice, puis vérifiez les statistiques et le projet confidentiel avec
Bob. Recommencez ces accès avec Alice et sans compte. Utilisez curl ou
l'onglet Réseau pour contrôler les statuts précis.

\subsection{10. Aide et bilan}\label{aide-et-bilan}

{\def\LTcaptype{none} % do not increment counter
\begin{longtable}[]{@{}
  >{\raggedright\arraybackslash}p{(\linewidth - 2\tabcolsep) * \real{0.5000}}
  >{\raggedright\arraybackslash}p{(\linewidth - 2\tabcolsep) * \real{0.5000}}@{}}
\toprule\noalign{}
\begin{minipage}[b]{\linewidth}\raggedright
Symptôme
\end{minipage} & \begin{minipage}[b]{\linewidth}\raggedright
Vérification
\end{minipage} \\
\midrule\noalign{}
\endhead
\bottomrule\noalign{}
\endlastfoot
API inaccessible & Spring doit écouter sur le port 8080 \\
\texttt{401} & Identifiants absents ou invalides ; vérifier la branche
et le compte disponible \\
Mutation \texttt{403} en étapes 2 et 3 & CSRF par défaut ; comportement
attendu sans jeton \\
Mutation \texttt{403} après l'exercice 4 & Vérifier le TODO 4.6 et le
redémarrage de Spring \\
Statistiques \texttt{403} avec Alice après l'exercice 4 & Comportement
attendu : rôle \texttt{ADMIN} requis \\
Note \texttt{404} en branches 1 à 4 avec un compte autorisé & Route
absente, fournie seulement à partir de la branche 5 \\
Alice lit la note après l'exercice 5 & Vérifier \texttt{@PreAuthorize}
et \texttt{@EnableMethodSecurity} \\
Tests ignorés & Retirer \texttt{@Disabled} seulement après avoir terminé
l'exercice correspondant \\
Ancien comportement après un changement de branche & Arrêter puis
relancer Spring et Vite \\
\end{longtable}
}

\texttt{401} indique une authentification absente ou invalide.
\texttt{403} indique un refus de sécurité : rôle insuffisant dans le
parcours final, ou jeton CSRF absent aux étapes précédentes.
\texttt{404} indique une route ou une tâche absente, lorsque la sécurité
autorise l'appel.

\textbf{Bilan à rendre :} fiche d'analyse, configurations des exercices,
résultats des tests et observations dans Planora. Expliquez la
distinction entre authentification et autorisation, l'intérêt de BCrypt,
l'ordre des règles HTTP et l'activation de la sécurité des méthodes.

HTTP Basic encode les identifiants en Base64 ; il ne les chiffre pas.
Les comptes et données de ce TP sont fictifs et destinés au local. Cette
démonstration n'est pas une configuration de production.

\textbf{Ressources :} \texttt{ENONCE.md} pour l'exercice courant ;
\texttt{docs/corrections/} pour les corrigés précédents ;
\texttt{docs/ARCHITECTURE.md}, \texttt{docs/SCENARIOS\_DE\_TEST.md} et
\texttt{docs/GUIDE\_HTTP.md} pour les détails.

\end{document}
